\chapter{Planck-Einheiten und universelle Konstanten in der fraktalen T0-Geometrie}


	
	\subsection*{Kurze Einführung}
	
	Dieses Kapitel leitet die Planck-Einheiten und alle fundamentalen Konstanten aus dem einzigen Parameter \(\xi\) ab.
	
	\subsection*{Mathematische Grundlage}
	
	Planck-Einheiten gelten als natürliche Skalen, bleiben aber im Standardmodell willkürlich. In der FFGFT emergieren sie aus der fraktalen Vakuumstruktur mit \(\xi = \frac{4}{3} \times 10^{-4}\).
	
	\subsection*{Fundamentale T0-Länge \(L_0\)}
	
	Die minimale Sub-Planck-Längenskala ist die fundamentale T0-Länge:
	
	\begin{equation}
		\boxed{L_0 = \xi \cdot l_P = \frac{4}{3} \times 10^{-4} \times 1{,}616 \times 10^{-35} \text{ m} = 2{,}155 \times 10^{-39} \text{ m}}
	\end{equation}
	
	\(L_0\) ist die Skala der Raumzeit-Granulation -- der absolute Boden der Physik, unter dem keine definierte Raumzeitgeometrie existiert. Sie ist etwa 7500 mal kleiner als die Planck-Länge.
	
	\subsection*{Planck-Länge als emergente Größe}
	
	Die Planck-Länge ist in der FFGFT eine abgeleitete Größe:
	
	\begin{equation}
		\ell_P = \frac{L_0}{\xi} = L_0 \cdot \xi^{-1}.
	\end{equation}
	
	Die Ableitungskette lautet: \(\xi \to G \to l_P\), nicht umgekehrt. Die Planck-Skala markiert den Übergang von der fraktalen Sub-Planck-Struktur zur klassischen Physik.
	
	\textbf{Einheitenprüfung:}
	\[
	[\ell_P] = \si{\meter}.
	\]
	
	\subsection*{Planck-Masse}
	
	Planck-Masse:
	
	\begin{equation}
		m_P = \rho_0 \cdot L_0^3 \cdot \xi^{-3}.
	\end{equation}
	
	Die Dichte \(\rho_0\) im Volumen \(L_0^3\), verstärkt durch \(\xi^{-3}\).
	
	\subsection*{Planck-Zeit}
	
	Planck-Zeit:
	
	\begin{equation}
		t_P = \frac{\ell_P}{c} = \frac{L_0}{c \cdot \xi}.
	\end{equation}
	
	Direkt aus Länge und Lichtgeschwindigkeit. Analog definiert \(T_0 = \xi \cdot t_P\) die minimale Zeitskala.
	
	\subsection*{Emergenz von \(G\), \(\hbar\), \(c\)}
	
	Gravitationskonstante:
	
	\begin{equation}
		G = \frac{\hbar c}{m_P^2} \cdot \xi^2.
	\end{equation}
	
	Schwäche durch \(\xi^2\).
	
	Alle Konstanten reduzieren auf \(\xi\), \(L_0\), \(\rho_0\), \(c\).
	
	\subsection*{Vergleich Standard -- FFGFT}
	
	\begin{center}
		\begin{tabular}{p{0.45\textwidth}p{0.45\textwidth}}
			\textbf{Standard} & \textbf{FFGFT (T0)} \\
			\hline
			Planck-Einheiten willkürlich & Emergent aus \(\xi\) \\
			19 freie Konstanten & Reduziert auf \(\xi\) \\
			Keine Hierarchie & Geometrisch erklärt \\
			Ad-hoc & Parameterfrei \\
		\end{tabular}
	\end{center}
	
	\subsection*{Schlussfolgerung}
	
	Die FFGFT leitet Planck-Einheiten und alle universellen Konstanten aus der fraktalen Skala \(\xi\) ab. Die fundamentale T0-Länge \(L_0 = \xi \cdot l_P \approx 2{,}155 \times 10^{-39}\,\)m ist die Sub-Planck-Skala der Raumzeit-Granulation. Die Hierarchieprobleme verschwinden -- alles ist geometrische Konsequenz eines einzigen Parameters.



